\chapter{Conclusion}\label{ch:conclusion}

This thesis examined how intrinsically interpretable credit-scoring models, particularly the Explainable Boosting Machine (EBM), compare with post-hoc-explained boosted-tree ensembles in predictive and profit-based performance and what the observed differences imply for model choice in the EU regulatory context. Under the empirical conditions studied, choosing the EBM did not entail a consistent performance penalty relative to XGBoost and CatBoost. A predictive and economic ensemble advantage emerged on Lending Club, but this pattern did not recur across the other datasets. The findings establish neither model-class equivalence nor a universally preferable architecture. They instead support including functionally flexible glass-box models in the initial empirical model-selection process rather than treating them solely as transparency-oriented alternatives. In the EU regulatory context, neither architecture is compliant or non-compliant by construction; conformity must be assessed for the complete deployed system.

The contrast between logistic regression and the EBM shows why the performance of a transparent linear baseline cannot be generalised to intrinsically interpretable models as a broader class. By combining direct inspectability with greater functional flexibility, the EBM remained more competitive with the boosted ensembles than logistic regression. Expected Maximum Profit largely reinforced the ordering indicated by discrimination while expressing the observed differences in economic terms under the stated assumptions. The explanation analysis found no consistent case-sampling stability advantage for the intrinsically interpretable pipelines. Although the EBM and the ensembles frequently prioritised similar source features, the analysis did not test whether these rankings accurately represented the models' predictive behaviour or were useful to credit decision-makers or affected applicants. The observed performance gaps also changed with training-set size, while the Lending Club ordering recurred under chronological evaluation. Within these sensitivity analyses, the precise performance gaps varied more than the model ordering.

These conclusions remain bounded by historical public datasets containing only booked loans, without reject inference or group-fairness evaluation. The principal comparisons rely on one train--test partition per dataset and a single temporal design, while EMP remains conditional on assumed economic parameters and ROC-AUC-based model selection. Moreover, the explanation analysis is restricted to the reproducibility and similarity of global feature rankings. These limitations restrict direct generalisation to contemporary institutional credit-scoring systems.

Future research should therefore evaluate these pipelines on current institutional data covering broader applicant populations, using product-specific economic assumptions and explicit fairness analyses. Repeated end-to-end model selection, EMP-based or multi-objective tuning and rolling-origin validation could determine whether the observed performance patterns persist when selection and temporal uncertainty are propagated more fully. Additional tests should examine whether the resulting explanations accurately represent model behaviour, whether they are understandable and useful to their intended audiences, and whether they can support the provision of intelligible individual explanations where legally required. Controlled organisational studies could further investigate whether direct inspectability reduces the effort required to document, validate and communicate a credit-scoring system under the applicable regulatory requirements.


